\documentclass[10pt,a4paper]{amsart}
\usepackage{amsmath,amssymb,amsthm,graphicx}
% Retained minimal preamble for the Stacks Project chapter more-algebra.tex (More on Algebra).
% The chapter file pinned-original-source.tex carries only \input{preamble} on its first line, so no
% preamble text lives inside the chapter itself.  The Stacks Project uses one single shared preamble for
% every chapter; the snapshot stages it identically as sources/stacks-chow/preamble.tex and
% sources/stacks-etale-cohomology/preamble.tex.  The block below is copied mechanically, line for line,
% from sources/stacks-chow/preamble.tex lines 14-18 (the Xy-pic setup) and 154-217 (the native theorem
% declarations together with the mathematical macros).  Nothing was added, removed or redefined; the
% wrapper class and the amsmath/amssymb/amsthm/graphicx base packages are supplied by the assembler.
% Chapter: more-algebra.tex -- More on Algebra.
% For commutative diagrams we use Xy-pic
\usepackage[all]{xy}

% We use 2cell for 2-commutative diagrams.
\xyoption{2cell}

% Theorem environments.
%
\theoremstyle{plain}
\newtheorem{theorem}[subsection]{Theorem}
\newtheorem{proposition}[subsection]{Proposition}
\newtheorem{lemma}[subsection]{Lemma}

\theoremstyle{definition}
\newtheorem{definition}[subsection]{Definition}
\newtheorem{example}[subsection]{Example}
\newtheorem{exercise}[subsection]{Exercise}
\newtheorem{situation}[subsection]{Situation}

\theoremstyle{remark}
\newtheorem{remark}[subsection]{Remark}
\newtheorem{remarks}[subsection]{Remarks}

\numberwithin{equation}{subsection}

% Macros
%
\def\lim{\mathop{\mathrm{lim}}\nolimits}
\def\colim{\mathop{\mathrm{colim}}\nolimits}
\def\Spec{\mathop{\mathrm{Spec}}}
\def\Hom{\mathop{\mathrm{Hom}}\nolimits}
\def\Ext{\mathop{\mathrm{Ext}}\nolimits}
\def\SheafHom{\mathop{\mathcal{H}\!\mathit{om}}\nolimits}
\def\SheafExt{\mathop{\mathcal{E}\!\mathit{xt}}\nolimits}
\def\Sch{\mathit{Sch}}
\def\Mor{\mathop{\mathrm{Mor}}\nolimits}
\def\Ob{\mathop{\mathrm{Ob}}\nolimits}
\def\Sh{\mathop{\mathit{Sh}}\nolimits}
\def\NL{\mathop{N\!L}\nolimits}
\def\CH{\mathop{\mathrm{CH}}\nolimits}
\def\proetale{{pro\text{-}\acute{e}tale}}
\def\etale{{\acute{e}tale}}
\def\QCoh{\mathit{QCoh}}
\def\Ker{\mathop{\mathrm{Ker}}}
\def\Im{\mathop{\mathrm{Im}}}
\def\Coker{\mathop{\mathrm{Coker}}}
\def\Coim{\mathop{\mathrm{Coim}}}

% Boxtimes
%
\DeclareMathSymbol{\boxtimes}{\mathbin}{AMSa}{"02}

%
% Macros for moduli stacks/spaces
%
\def\QCohstack{\mathcal{QC}\!\mathit{oh}}
\def\Cohstack{\mathcal{C}\!\mathit{oh}}
\def\Spacesstack{\mathcal{S}\!\mathit{paces}}
\def\Quotfunctor{\mathrm{Quot}}
\def\Hilbfunctor{\mathrm{Hilb}}
\def\Curvesstack{\mathcal{C}\!\mathit{urves}}
\def\Polarizedstack{\mathcal{P}\!\mathit{olarized}}
\def\Complexesstack{\mathcal{C}\!\mathit{omplexes}}
% \Pic is the operator that assigns to X its picard group, usage \Pic(X)
% \Picardstack_{X/B} denotes the Picard stack of X over B
% \Picardfunctor_{X/B} denotes the Picard functor of X over B
\def\Pic{\mathop{\mathrm{Pic}}\nolimits}
\def\Picardstack{\mathcal{P}\!\mathit{ic}}
\def\Picardfunctor{\mathrm{Pic}}
\def\Deformationcategory{\mathcal{D}\!\mathit{ef}}

% --- appended 2026-10-12 (Chief): the project own definitions of the reference/slogan/history
% display environments, copied mechanically from the shared Stacks preamble lines 8-12
% (\usepackage{verbatim} plus the three \newenvironment lines, which the chapter preamble keeps
% verbatim from the project preamble). Nothing added, removed or redefined: these environments
% carry no mathematical content and the project itself typesets them as comments. Without these
% three lines a counted statement whose own text uses them fails with Environment slogan undefined.
\usepackage{verbatim}
\newenvironment{reference}{\comment}{\endcomment}
\newenvironment{slogan}{\comment}{\endcomment}
\newenvironment{history}{\comment}{\endcomment}
\usepackage{cleveref}
\crefname{originallemma}{original lemma}{original lemmas}
\Crefname{originallemma}{Original lemma}{Original lemmas}
\makeatletter
\newlabel{more-algebra-lemma-glue-pseudo-coherent}{{066D}{source}{}{sourcepointer-1}{}}
\newlabel{more-algebra-lemma-glue-pseudo-coherent@cref}{{[originallemma][0][]066D}{[1][0][]source}{}{}{}}
\makeatother
\crefname{originallemma}{original lemma}{original lemmas}
\Crefname{originallemma}{Original lemma}{Original lemmas}
\makeatletter
\newlabel{algebra-lemma-cover}{{00EO}{source}{}{sourcepointer-2}{}}
\newlabel{algebra-lemma-cover@cref}{{[originallemma][0][]00EO}{[1][0][]source}{}{}{}}
\makeatother
\crefname{originallemma}{original lemma}{original lemmas}
\Crefname{originallemma}{Original lemma}{Original lemmas}
\makeatletter
\newlabel{lemma-recognize-pseudo-coherent}{{064W}{source}{}{sourcepointer-3}{}}
\newlabel{lemma-recognize-pseudo-coherent@cref}{{[originallemma][0][]064W}{[1][0][]source}{}{}{}}
\makeatother
\crefname{originallemma}{original lemma}{original lemmas}
\Crefname{originallemma}{Original lemma}{Original lemmas}
\makeatletter
\newlabel{lemma-finite-cohomology}{{064S}{source}{}{sourcepointer-4}{}}
\newlabel{lemma-finite-cohomology@cref}{{[originallemma][0][]064S}{[1][0][]source}{}{}{}}
\makeatother
\crefname{originallemma}{original lemma}{original lemmas}
\Crefname{originallemma}{Original lemma}{Original lemmas}
\makeatletter
\newlabel{lemma-cone-pseudo-coherent}{{064R}{source}{}{sourcepointer-5}{}}
\newlabel{lemma-cone-pseudo-coherent@cref}{{[originallemma][0][]064R}{[1][0][]source}{}{}{}}
\makeatother
\title{Pseudo-coherence is local for a finite cover by principal open sets}
\author{The Stacks Project Authors}
\date{stacks-project:more-algebra.tex:local-snapshot-20261012; GNU FDL 1.2 or later}
\begin{document}
\maketitle
\setcounter{subsection}{1}
\noindent\textit{Editorial note.} All delivered mathematical text is the author's own, copied from the pinned Stacks Project chapter file more-algebra.tex; no statement or proof step was rewritten, reordered or removed.\par\medskip
\noindent\textit{Editorial note.} The chapter file pinned-original-source.tex carries only \textbackslash{}input\{preamble\} on its first line, so no preamble text lives inside the chapter. The retained file sources/stacks-more-algebra/retained-preamble.tex is a mechanical, line-for-line copy of the project's single shared preamble as staged in the snapshot (Xy-pic setup plus theorem declarations and mathematical macros). Nothing was added, removed or redefined, and the excerpt serves only as preamble, never as a span coordinate.\par\medskip
\noindent\textit{Editorial note.} Both counted spans start at the \textbackslash{}begin line of their environment, so the statement and the author's proof are delivered whole; surrounding prose of the section is deliberately outside the counted spans and is carried as standard prerequisites instead.\par\medskip
\noindent\textit{Editorial note.} The citation fields (source\_id, source\_version, authors, source\_url, license.name, source\_title) are identical to those of every other card this lane has produced from this chapter, so the batch selector treats them as one source.\par\medskip
\noindent\textit{Allowed standard prerequisites (editorial).}\begin{itemize}
\item Finiteness of a module can be checked on a finite cover by principal open sets. Algebra, Lemma \textbackslash{}ref\{algebra-lemma-cover\}, cited at proof lines 16421 and 16431. Turn finiteness of the localised cohomology modules into finiteness over the ring, and vanish the cohomology above the top degree.
\item The m-pseudo-coherent complexes are exactly those with vanishing cohomology above m and finitely generated cohomology in degree m. Lemma \textbackslash{}ref\{lemma-recognize-pseudo-coherent\}, cited at proof lines 16424-16425. Gives the base case of the induction.
\item Finiteness of cohomology in the top nonvanishing degree characterises pseudo-coherence one degree at a time. Lemma \textbackslash{}ref\{lemma-finite-cohomology\}, cited at proof lines 16430-16431. Provides finiteness of the top cohomology module after localisation.
\item Cone lemma for pseudo-coherence in a distinguished triangle. Lemma \textbackslash{}ref\{lemma-cone-pseudo-coherent\}, cited at proof lines 16439-16440 and 16443-16444. Keeps the cone m-pseudo-coherent after localisation and finishes the induction.
\end{itemize}\medskip
\ifdefined\hypertarget\hypertarget{sourcepointer-1}{}\fi
\noindent\textit{Source reference.} Original lemma 066D: focused-seeds/AGP/sources/stacks-snapshot-20261011/tags.txt, source lines 7932--7932. The lemma's own canonical tag, quoted so that readers can resolve the label inside the delivered text; it points at the counted statement itself and carries no proof content.\par\medskip
\ifdefined\hypertarget\hypertarget{sourcepointer-2}{}\fi
\noindent\textit{Source reference.} Original lemma 00EO: focused-seeds/AGP/sources/stacks-snapshot-20261011/tags.txt, source lines 542--542. Navigation pointer to a statement the counted proof refers to but does not contain; nothing load-bearing is delegated to it.\par\medskip
\ifdefined\hypertarget\hypertarget{sourcepointer-3}{}\fi
\noindent\textit{Source reference.} Original lemma 064W: focused-seeds/AGP/sources/stacks-snapshot-20261011/tags.txt, source lines 7880--7880. Navigation pointer to a statement the counted proof refers to but does not contain; nothing load-bearing is delegated to it.\par\medskip
\ifdefined\hypertarget\hypertarget{sourcepointer-4}{}\fi
\noindent\textit{Source reference.} Original lemma 064S: focused-seeds/AGP/sources/stacks-snapshot-20261011/tags.txt, source lines 7876--7876. Navigation pointer to a statement the counted proof refers to but does not contain; nothing load-bearing is delegated to it.\par\medskip
\ifdefined\hypertarget\hypertarget{sourcepointer-5}{}\fi
\noindent\textit{Source reference.} Original lemma 064R: focused-seeds/AGP/sources/stacks-snapshot-20261011/tags.txt, source lines 7875--7875. Navigation pointer to a statement the counted proof refers to but does not contain; nothing load-bearing is delegated to it.\par\medskip
\begin{lemma}
\label{lemma-glue-pseudo-coherent}
Let $R$ be a ring. Let $f_1, \ldots, f_r \in R$ be elements which
generate the unit ideal. Let $m \in \mathbf{Z}$. Let $K^\bullet$
be a complex of $R$-modules. If for each $i$ the complex
$K^\bullet \otimes_R R_{f_i}$ is $m$-pseudo-coherent
(resp.\ pseudo-coherent), then $K^\bullet$ is $m$-pseudo-coherent
(resp.\ pseudo-coherent).
\end{lemma}

\begin{proof}
We will use without further mention that $- \otimes_R R_{f_i}$ is
an exact functor and that therefore
$$
H^i(K^\bullet)_{f_i} =
H^i(K^\bullet) \otimes_R R_{f_i} = H^i(K^\bullet \otimes_R R_{f_i}).
$$
Assume $K^\bullet \otimes_R R_{f_i}$ is $m$-pseudo-coherent
for $i = 1, \ldots, r$. Let $n \in \mathbf{Z}$ be the largest
integer such that $H^n(K^\bullet \otimes_R R_{f_i})$ is nonzero
for some $i$. This implies in particular that $H^i(K^\bullet) = 0$
for $i > n$ (and that $H^n(K^\bullet) \not = 0$) see
Algebra, Lemma \ref{algebra-lemma-cover}.
We will prove the lemma by induction on $n - m$.
If $n < m$, then the lemma is true by
Lemma \ref{lemma-recognize-pseudo-coherent}.
If $n \geq m$, then $H^n(K^\bullet)_{f_i}$ is a finite $R_{f_i}$-module
for each $i$, see
Lemma \ref{lemma-finite-cohomology}.
Hence $H^n(K^\bullet)$ is a finite $R$-module, see
Algebra, Lemma \ref{algebra-lemma-cover}.
Choose a finite free $R$-module $E$ and a surjection $E \to H^n(K^\bullet)$.
As $E$ is projective we can lift this to a map of complexes
$\alpha : E[-n] \to K^\bullet$. Then the cone $C(\alpha)^\bullet$ has
vanishing cohomology in degrees $\geq n$. On the other hand, the
complexes $C(\alpha)^\bullet \otimes_R R_{f_i}$ are $m$-pseudo-coherent
for each $i$, see
Lemma \ref{lemma-cone-pseudo-coherent}.
Hence by induction we see that $C(\alpha)^\bullet$ is $m$-pseudo-coherent
as a complex of $R$-modules. Applying
Lemma \ref{lemma-cone-pseudo-coherent}
once more we conclude.
\end{proof}

\end{document}
